% Generated draft: requires mathematical and source review.
\documentclass[11pt]{article}
\usepackage[margin=25mm]{geometry}
\usepackage{fontspec}
\setmainfont{Latin Modern Roman}
\usepackage{amsmath,amssymb,mathtools,unicode-math}
\setmathfont{Latin Modern Math}
\usepackage{enumitem,needspace,xurl,hyperref}
\setlength{\parindent}{0pt}
\setlength{\parskip}{5pt}
\setlength{\emergencystretch}{3em}
\newcommand{\N}{\mathbb{N}}
\newcommand{\Z}{\mathbb{Z}}
\newcommand{\Q}{\mathbb{Q}}
\newcommand{\R}{\mathbb{R}}
\newcommand{\C}{\mathbb{C}}
\newcommand{\F}{\mathbb{F}}
\newcommand{\A}{\mathbb{A}}
\newcommand{\PP}{\mathbb{P}}
\newcommand{\E}{\mathbb{E}}
\newcommand{\eps}{\varepsilon}
\newcommand{\dd}{\,\mathrm{d}}
\begin{document}
\section*{UnsolvedMath ID: 30002039}
\textbf{Source problem code:} catalogue entry
\subsection{Definitions and mathematical statement}
A deformation of projective surfaces is understood here as a flat projective family \(\pi\colon\mathcal X\to B\) with a distinguished point \(0\in B\) and special fiber \(\mathcal X_0=X\). The supplied record does not specify the ground field, the base \(B\), or additional hypotheses on the total space. Nonisotriviality means that the isomorphism class of the fibers varies; the precise generic or local formulation is not specified.

In the usual characteristic-zero surface terminology, a \(T\)-singularity is a quotient surface singularity admitting a \(\Q\)-Gorenstein smoothing. The \(\Q\)-Gorenstein condition involves the relative canonical divisor being \(\Q\)-Cartier. For families, conventions may additionally require suitable reflexive powers of the relative dualizing sheaf to commute with base change. The record does not specify these conventions or the characteristic, so this terminology does not supply missing hypotheses silently.

For an irreducible projective surface \(Y\), a smooth minimal model means a smooth projective surface \(M\) birational to \(Y\), with no exceptional curve of the first kind, that is, no smooth rational curve of self-intersection \(-1\). When \(Y\) is singular, this concerns a smooth model of its birational class, obtainable by resolving singularities and then contracting exceptional curves where those operations are available. It does not mean that \(Y\) itself is smooth, and uniqueness of \(M\) is not assumed.

\textbf{Classification question.} Given a nonisotrivial \(\Q\)-Gorenstein deformation \(\pi\colon\mathcal X\to B\) of projective surfaces whose special fiber \(X\) is irreducible and normal and has \(T\)-singularities, determine the smooth minimal models of the singular and nonsingular surfaces represented in the moduli space near the stable limit of \(X\).

This preserves the supplied question, but its classification target remains underdetermined: the record does not identify the moduli space, the meaning of its neighborhood, the construction or existence hypotheses for the stable limit, or whether the surfaces to be classified range over the given deformation or over a larger moduli neighborhood. If ``stable'' refers to KSBA stability, the relevant surfaces are semi-log-canonical with ample canonical divisor, but the source record does not explicitly impose this interpretation or the hypotheses needed to obtain such a limit. It also does not specify whether models are to be classified up to isomorphism or with additional markings. These qualifications require resolution before the question becomes a fully specified classification problem.
\subsection{Short English statement}
For a nonisotrivial \(\Q\)-Gorenstein family of projective surfaces with irreducible normal special fiber having \(T\)-singularities, describe the smooth minimal models of nearby singular and smooth surfaces around its stable limit. The field, moduli space, stable-limit construction, neighborhood, and classification equivalence are unspecified, so this is a qualified classification question requiring mathematical review.
\subsection{Sources}
\begin{itemize}
\item Supplied local UnsolvedMath registry record, permanent ID \(30002039\), title \emph{Minimal Models Near Stable Surface Limits}, source problem code \texttt{catalogue entry}, in the Oberwolfach Reports Open Problems collection. The supplied proposed year is 2012; no individual proposer is identified.
\item Supplied report citation: \emph{Toric Geometry (2013)}, Oberwolfach Reports, DOI \texttt{10.4171/owr/2012/21}, pp.\ 1231--1306. \url{https://doi.org/10.4171/owr/2012/21}. The record labels the statement \texttt{corrected\_verified} and reports that the official report was checked, the complete question preserved, and deformation and singularity notation normalized. This is supplied verification metadata, not an independent verification by this draft.
\item Hacking--Tevelev--Urz\'ua, \emph{Flipping surfaces}. \url{https://arxiv.org/abs/1310.1580}.
\item Tevelev--Urz\'ua, \emph{Categorical aspects of the KSB correspondence}. \url{https://arxiv.org/abs/2204.13225}.
\item Literature assessment dated exactly \texttt{2026-08-21}: the supplied curation labels the problem partially solved and states that Hacking--Tevelev--Urz\'ua give an algorithm and classify substantial classes of flips and antiflips, that later work treats \(N\)-resolutions, and that the broad question concerning all smooth minimal models near an arbitrary stable limit is not completely classified. This is supplied dated curation derived from the OWR material, not an independently verified literature assessment; its status and cited claims require independent checking.
\end{itemize}
\end{document}
